% Options for packages loaded elsewhere
\PassOptionsToPackage{unicode}{hyperref}
\PassOptionsToPackage{hyphens}{url}
\documentclass[
]{article}
\usepackage{xcolor}
\usepackage[margin=1in]{geometry}
\usepackage{amsmath,amssymb}
\setcounter{secnumdepth}{-\maxdimen} % remove section numbering
\usepackage{iftex}
\ifPDFTeX
  \usepackage[T1]{fontenc}
  \usepackage[utf8]{inputenc}
  \usepackage{textcomp} % provide euro and other symbols
\else % if luatex or xetex
  \usepackage{unicode-math} % this also loads fontspec
  \defaultfontfeatures{Scale=MatchLowercase}
  \defaultfontfeatures[\rmfamily]{Ligatures=TeX,Scale=1}
\fi
\usepackage{lmodern}
\ifPDFTeX\else
  % xetex/luatex font selection
\fi
% Use upquote if available, for straight quotes in verbatim environments
\IfFileExists{upquote.sty}{\usepackage{upquote}}{}
\IfFileExists{microtype.sty}{% use microtype if available
  \usepackage[]{microtype}
  \UseMicrotypeSet[protrusion]{basicmath} % disable protrusion for tt fonts
}{}
\makeatletter
\@ifundefined{KOMAClassName}{% if non-KOMA class
  \IfFileExists{parskip.sty}{%
    \usepackage{parskip}
  }{% else
    \setlength{\parindent}{0pt}
    \setlength{\parskip}{6pt plus 2pt minus 1pt}}
}{% if KOMA class
  \KOMAoptions{parskip=half}}
\makeatother
\usepackage{graphicx}
\makeatletter
\newsavebox\pandoc@box
\newcommand*\pandocbounded[1]{% scales image to fit in text height/width
  \sbox\pandoc@box{#1}%
  \Gscale@div\@tempa{\textheight}{\dimexpr\ht\pandoc@box+\dp\pandoc@box\relax}%
  \Gscale@div\@tempb{\linewidth}{\wd\pandoc@box}%
  \ifdim\@tempb\p@<\@tempa\p@\let\@tempa\@tempb\fi% select the smaller of both
  \ifdim\@tempa\p@<\p@\scalebox{\@tempa}{\usebox\pandoc@box}%
  \else\usebox{\pandoc@box}%
  \fi%
}
% Set default figure placement to htbp
\def\fps@figure{htbp}
\makeatother
\setlength{\emergencystretch}{3em} % prevent overfull lines
\providecommand{\tightlist}{%
  \setlength{\itemsep}{0pt}\setlength{\parskip}{0pt}}
\usepackage{bookmark}
\IfFileExists{xurl.sty}{\usepackage{xurl}}{} % add URL line breaks if available
\urlstyle{same}
\hypersetup{
  pdftitle={test},
  hidelinks,
  pdfcreator={LaTeX via pandoc}}

\title{test}
\author{}
\date{\vspace{-2.5em}}

\begin{document}
\maketitle

「クチャクチャ」という音（いわゆるクチャラー）は、自分では気づきにくいものですが、一度気になると改善したい大切なマナーですよね。その心掛け、素晴らしいと思います！

音が出てしまう原因の多くは、\textbf{「口の中に隙間があること」}と\textbf{「呼吸のタイミング」}にあります。今日から意識できる具体的なコツをまとめました。

\begin{center}\rule{0.5\linewidth}{0.5pt}\end{center}

\subsection{1.
物理的に音を遮断する「鉄則」}\label{ux7269ux7406ux7684ux306bux97f3ux3092ux906eux65adux3059ux308bux9244ux5247}

まずは、音を物理的に鳴らさないための基本的な動作を見直しましょう。

\begin{itemize}
\tightlist
\item
  \textbf{唇をしっかり閉じる}
  噛んでいる間、上下の唇をぴたっと密着させることを最優先にしてください。少しでも隙間があると、そこから空気が漏れて音が発生します。
\item
  \textbf{一口の量を減らす}
  口いっぱいに食べ物を入れると、唇を閉じるのが難しくなり、咀嚼（そしゃく）の際に音が漏れやすくなります。「今の半分」を意識して口に運んでみてください。
\item
  \textbf{「奥歯」で噛む}
  前歯の方で噛むと唇が開きやすくなります。食べ物を口の奥へ送り、奥歯で静かに噛むようにしましょう。
\end{itemize}

\begin{center}\rule{0.5\linewidth}{0.5pt}\end{center}

\subsection{2.
呼吸と姿勢を整える}\label{ux547cux5438ux3068ux59ffux52e2ux3092ux6574ux3048ux308b}

意外な盲点が「鼻」と「姿勢」です。

\begin{itemize}
\tightlist
\item
  \textbf{鼻呼吸を徹底する}
  口を閉じて噛むためには、鼻で息ができなければなりません。もし鼻が詰まっている場合は、どうしても口で息をしながら噛むことになり、音が鳴ってしまいます。
\item
  \textbf{姿勢を正す}
  猫背になると喉が圧迫され、呼吸が苦しくなって口が開きやすくなります。背筋を伸ばし、食べ物を迎えに行くのではなく、食器を口元に運ぶようにしましょう。
\end{itemize}

\begin{center}\rule{0.5\linewidth}{0.5pt}\end{center}

\subsection{3.
意識を変えるトレーニング}\label{ux610fux8b58ux3092ux5909ux3048ux308bux30c8ux30ecux30fcux30cbux30f3ux30b0}

無意識の習慣を上書きするためのステップです。

\begin{itemize}
\tightlist
\item
  \textbf{自分の音を「録音」してみる}
  一度、スマホで自分の食事風景を録音・録画してみてください。客観的に自分の音を聞くことで、「どのタイミングで音が鳴るか」が明確になり、意識が劇的に変わります。
\item
  \textbf{飲み込んでから喋る}
  口の中に食べ物がある状態で話そうとすると、必ず唇が開いて音が漏れます。「完全に飲み込むまで口を開けない」を徹底しましょう。
\end{itemize}

\begin{quote}
\textbf{💡 ワンポイント・アドバイス}
慣れるまでは、静かな場所よりも少しガヤガヤした場所や、テレビがついている環境で練習すると、プレッシャーを感じすぎずに取り組めますよ。
\end{quote}

\begin{center}\rule{0.5\linewidth}{0.5pt}\end{center}

\subsubsection{次に私ができること}\label{ux6b21ux306bux79c1ux304cux3067ux304dux308bux3053ux3068}

まずは次の食事で「唇を閉じる」ことだけを意識してみてください。もし「どうしても鼻が詰まっていて息がしづらい」や「特定の食べ物（麺類など）だけ難しい」といった具体的なお悩みがあれば、さらに踏み込んだ対策を提案できます。

\textbf{次は、麺類を静かに食べるコツや、外食時のマナーについて詳しくお伝えしましょうか？}

\end{document}
